% ===========================================================================
% ACT I -- Opening and Motivation  (Ch. 1)   Slides 1-8
% Budget: 5 minutes.
% ===========================================================================

% --- 1. Title ------------------------------------------------------------
{%
\setbeamertemplate{footline}{}
\begin{frame}[plain]
  \vspace{0.2em}
  \begin{center}
    {\color{uniSlate}\footnotesize\textsc{Colorado State University Pueblo}
     \\[0.15em] \scriptsize School of Engineering}

    \vspace{0.9em}
    \color{uniCrimson}\rule{0.55\textwidth}{2pt}

    \vspace{0.8em}
    {\color{uniNavy}\LARGE\bfseries A Graph Neural Surrogate for\\[0.25em]
     Longitudinal Train Dynamics\par}

    \vspace{0.5em}
    {\color{uniSlate}\normalsize Toward Fuel-Optimal Control of Freight Consists\par}

    \vspace{0.8em}
    \color{uniCrimson}\rule{0.55\textwidth}{2pt}

    \vspace{0.8em}
    {\color{uniInk}\large George James\par}
    \vspace{0.2em}
    {\color{uniSlate}\footnotesize Master's Thesis Defense
      \quad\textbullet\quad \today\par}

    \vspace{0.6em}
    \scalebox{0.80}{\consistBand}
  \end{center}
\end{frame}
}

% --- 2. What LTD are -----------------------------------------------------
\begin{frame}{Background: What Longitudinal Train Dynamics Are}
  \begin{columns}[T,onlytextwidth]
    \begin{column}{0.42\textwidth}
      \begin{itemize}\setlength\itemsep{0.6em}
        \item A freight train is a \textbf{distributed} mechanical system:
              tens to hundreds of coupled masses.
        \item Each vehicle feels the grade, curvature and resistance
              at its \textbf{own} position on the route.
        \item Couplers carry \alert{slack}: free travel before any force
              is transmitted.
        \item Quantities that matter in operation: \textbf{fuel},
              \textbf{arrival time}, \textbf{peak coupler force}.
      \end{itemize}
    \end{column}
    \begin{column}{0.56\textwidth}
      \centering
      \vspace{0.6em}
      \scalebox{0.70}{\trainFBD}

      \vspace{1.8em}
      \begin{beamercolorbox}[wd=0.92\textwidth,sep=6pt,center]{block body}
        $m_i\dot v_i \;=\; F_{\mathrm{ext},i} + F_{i-1} - F_i$
      \end{beamercolorbox}
    \end{column}
  \end{columns}
\end{frame}

% --- 3. Slack ------------------------------------------------------------
\begin{frame}{Coupler Slack Makes the Problem Hard}
  \centering
  \scalebox{0.74}{%
    \only<1>{\slackScene{1}}%
    \only<2>{\slackScene{2}}%
    \only<3>{\slackScene{3}}%
  }

  \vspace{0.3em}
  \only<3>{%
    \begin{minipage}{0.86\textwidth}\small
      When slack closes suddenly, the force spike can break a coupler. The
      same \alert{abrupt change in force} makes the numerical solver take very
      small time steps, which is why the simulation is slow.
    \end{minipage}}
\end{frame}

% --- 4. Published cost ---------------------------------------------------
\begin{frame}{High-Fidelity Simulation Is Too Slow to Use in a Loop}
  \vspace{0.3em}
  \centering
  \begin{tabular}{@{}llr@{}}
    \toprule
    \textbf{Task} & \textbf{Reported Cost} & \textbf{Source} \\
    \midrule
    Whole-trip LTD optimization, sequential  & $\sim$18 months
      & {\footnotesize Wu et al.\ (2016)} \\
    \quad the same study, run in parallel    & 11 days
      & {\footnotesize Wu et al.\ (2016)} \\[0.35em]
    Multibody vehicle--track simulation, 5\,km & $\sim$30 minutes
      & {\footnotesize Zhou et al.\ (2023)} \\
    \quad \textbf{learned surrogate, same 5\,km}
      & \textcolor{uniCrimson}{\textbf{$\sim$8 seconds}}
      & {\footnotesize Zhou et al.\ (2023)} \\
    \bottomrule
  \end{tabular}

  \vspace{1.0em}
  \begin{minipage}{0.88\textwidth}\small
    At this cost, any task that needs thousands of simulations is impractical:
    optimizing a driving plan, testing many scenarios, or running a controller.
  \end{minipage}

  \vspace{0.7em}
  \srcline{Figures quoted as reported; not normalized to a common hardware baseline.}
\end{frame}

% --- 5. Our own benchmark ------------------------------------------------
\begin{frame}{Our Reference Simulator Has the Same Problem}
  \vspace{-0.2em}
  \begin{columns}[T,onlytextwidth]
    \begin{column}{0.58\textwidth}
      \benchRegimeChart
      \vspace{-0.2em}
      \centering\benchLegend
    \end{column}
    \begin{column}{0.40\textwidth}
      \vspace{0.2em}
      {\small Runs completed within the 30-second budget, by consist size:}
      \vspace{0.5em}

      \begin{center}
        \stattile{16 / 22}{$N = 11$ \quad median 4.0\,s}
        \vspace{0.45em}
        \stattile{4 / 22}{$N = 60$ \quad median 28.1\,s}
        \vspace{0.45em}
        \stattile{\textcolor{uniCrimson}{0 / 22}}{$N = 130$ \quad median $>$ 30\,s}
      \end{center}
    \end{column}
  \end{columns}

  \vspace{0.1em}
  \begin{minipage}{0.95\textwidth}\scriptsize
    The slowest runs are the ones with frequent slack events, because each event
    makes the solver shrink its step. These runs also produce the largest coupler
    forces, which are the forces a safety limit is meant to control.
  \end{minipage}
\end{frame}

% --- 6. Surrogate modelling ----------------------------------------------
\begin{frame}{Surrogate Modeling}
  \vspace{0.8em}
  \centering
  \scalebox{0.86}{\surrogateFlow}

  \vspace{1.4em}
  \begin{minipage}{0.90\textwidth}\small
    A surrogate is a fast model trained on simulator output. It gives up a small,
    measured amount of accuracy for a large gain in speed. It \textbf{does not
    replace} the simulator: the simulator produces the training data and remains
    the reference the surrogate is checked against.
  \end{minipage}
\end{frame}

% --- 7. Question and contributions ---------------------------------------
\begin{frame}{Thesis Question and Contributions}
  \vspace{0.2em}
  \begin{beamercolorbox}[wd=\textwidth,sep=8pt]{block body}
    \centering\itshape
    Can a train-agnostic, differentiable surrogate of longitudinal train
    dynamics be accurate enough, and fast enough, to sit inside a fuel-optimal
    control loop?
  \end{beamercolorbox}

  \vspace{0.8em}
  \begin{enumerate}\setlength\itemsep{0.45em}
    \item A \textbf{graph representation} of the train, with vehicles as nodes
          and couplers as edges, that works for any train length.
    \item A \textbf{GPU version of the reference simulator} that runs many
          scenarios in parallel, and two 10,000-scenario datasets built with it.
    \item A \textbf{graph network simulator} that advances the train by one
          learned 0.25\,s step, replacing 50 evaluations of the physics
          equations. It is compared against simple baselines, an MLP and a
          transformer.
    \item A \textbf{fuel-optimal control demonstration} (MPC) that uses the
          surrogate's gradients to improve the driving plan.
  \end{enumerate}
\end{frame}

% --- 8. Roadmap ----------------------------------------------------------
\begin{frame}{Roadmap}
  \vspace{0.4em}
  \begin{columns}[T,onlytextwidth]
    \begin{column}{0.46\textwidth}
      \centering
      \scalebox{0.88}{\actstrip{1}}
    \end{column}
    \begin{column}{0.50\textwidth}
      \vspace{1.2em}
      {\small
      The talk follows the thesis in order.

      \vspace{0.9em}
      \textbf{Acts 3--6} cover the formulation, the reference simulator and the
      dataset. These parts are complete and validated.

      \vspace{0.9em}
      \textbf{Acts 7--10} are the surrogate itself, its evaluation, and the
      control demonstration it enables.

      \vspace{0.9em}
      This strip appears at the start of each act to show where we are.}
    \end{column}
  \end{columns}
\end{frame}
